\documentclass[11pt]{article}
\usepackage[margin=0.85in]{geometry}
\usepackage{amsmath,booktabs,hyperref}
\hypersetup{colorlinks=true,urlcolor=blue}
\setlength{\parindent}{0pt}
\setlength{\parskip}{0.55em}
\title{Target Estimate Updates and Empirical Experiment Outcomes}
\author{Collision Avoidance Research}
\date{October 3, 2026}
\begin{document}
\maketitle

\textbf{Implemented behavior.} Every controller call uses the latest target
estimate. The estimated acceleration and sideslip are constant over that
call's prediction, while subsequent estimates may change:
\[
 \widehat A_{i|k}=\widehat A_k,\qquad
 \widehat\beta_{i|k}=\widehat\beta_k,\qquad
 (\widehat A_{k+1},\widehat\beta_{k+1})
 \text{ need not equal }(\widehat A_k,\widehat\beta_k).
\]
The previous input plan remains the normal initialization. Target prediction
updates do not require an enclosure certificate or an explicit controller
reset. Heading is unwrapped against the previous prediction to preserve
continuity. No assumption-status diagnostics are generated.

\section{Theory and executable experiments}
The conditional recursive-feasibility derivation in
\path{controller/OBSERVER_ROBUST_PCBF_THEORY.tex} remains a mathematical result
under its ideal premises. Experiments can exercise the same algorithm outside
those premises. An unavailable, stale, incomplete or misaligned enclosure no
longer throws an admission error. Usable current error bounds still tighten
the existing optimization; an unusable bound supplies no tightening and is
not treated as evidence of zero error. The existing terminal constraints,
single CLF, primary restoration and numerical model-agreement/damping logic
remain unchanged. Actual optimization failure still means that no control
can be issued.

The numerical interface uses one controller, without an executable backup or
an additional CLF. No observer equations, gains, uncertainty radii, collision
margin or physical target dynamics were retuned in this change.

\section{Independent truth and failure definition}
Previously the offline replay evaluated collision using the controller's
target epoch. That was adequate for fixed exact observations but would let
changed estimates alter the collision reference. The replay now advances
target truth solely from the scenario's initial condition. Synthetic
observations affect planning, never the physical target.

Failure is rectangle contact/overlap in the independent replay or failure to
return a finite solved control. A theoretical premise is not an experimental
failure criterion. Recovery is a separately observed dwell in the nominal
behavior tolerances. Reaching an observation cap without recovery is neither
a collision nor a claim of recovery. Replay/harness execution errors have a
separate inconclusive outcome. Historical reports are not rewritten.

The replay uses \texttt{ode45}, relative tolerance $10^{-11}$, absolute
tolerance $10^{-12}$ and 31 evaluations per 50 ms hold. It is finite sampled
evidence, not an interval collision certificate.

\section{Validation}
The full MATLAB suite passed \textbf{682/682}, with no failed or
incomplete tests. It covers updates of both estimated parameters within one
prediction, warm-start preservation, heading wrapping, unavailable/stale
enclosures, continued control after loss of a certificate, independent truth,
collision failure despite continuing solved controls, actual optimization
failure and observation caps without a false recovery claim. Existing
finite uncertainty constraints can still make the optimization infeasible;
the corresponding regression remains in place.

A separate probe consumed an actual \texttt{onlineNrmmTrackingRuntime}
publication initialized without an explicit error prior. Its ego speed was
10 m/s, the target relative position was $(20,5)$ m, and target velocity was
$(-12,0)$ m/s. The controller did not return a solved input: the PCBF stage returned no numerical solution. This is an actual
optimization outcome, not a rejection of a theoretical premise. The paired
study below uses synthetic target observations, so it does not resolve this
real-publication case or constitute a sensor/observer closed-loop validation.

\section{Paired closed-loop study}
Three 8 m/s collision-threat scenarios were run with exact observations and
with deterministic target-estimate perturbations. True target acceleration
and sideslip stayed constant in each scenario. The observation perturbations
were
\begin{align*}
 \Delta p&=0.03e^{-t/5}[\sin(3t),\cos(5t)]^\top\ \mathrm m,\\
 \Delta V&=0.02e^{-t/5}\sin(2t)\ \mathrm{m/s},\\
 \Delta A&=0.02e^{-t/5}\cos(4t)\ \mathrm{m/s^2},\\
 \Delta\beta&=0.0005e^{-t/5}\sin(6t)\ \mathrm{rad}.
\end{align*}
The perturbed publications carried unavailable certificates. Their body
heading was exact and the reported velocity direction was consistent with
the perturbed sideslip. Ego observations were exact. No random seed is used.
This isolates robustness to revised forecasts; it does not emulate a complete
observer or claim performance for all measurement errors.

The controller used an 8-step prefix and the existing completion horizon,
0.20 m collision-node margin, 10 m path-center deviation limit, and no road
boundary constraints. Each run continued until a one-second nominal-recovery
dwell or an actual failure, with 1200 holds as an observation cap. Tolerances
were 0.1 m lateral error, one degree heading error, 0.1 m/s speed error,
0.05 m/s lateral velocity and 0.01 rad/s yaw-rate error.

\begin{center}
\small
\begin{tabular}{llrrrr}
\toprule
Scenario & Observation & Holds & Min. gap (m) & Recovery (s) & Outcome\\
\midrule
Head-on & Exact & 204 & 1.3922 & 10.20 & recovered \\
Accelerating head-on & Exact & 205 & 1.0645 & 10.25 & recovered \\
Crossing & Exact & 254 & 0.1626 & 12.70 & recovered \\
Head-on & Perturbed & 204 & 1.3951 & 10.20 & recovered \\
Accelerating head-on & Perturbed & 205 & 1.0671 & 10.25 & recovered \\
Crossing & Perturbed & 253 & 0.1997 & 12.65 & recovered \\
\bottomrule
\end{tabular}
\end{center}

Timing samples in the JSON are controller-call wall times (including the
synthetic observation hook). Other processes were active; these measurements
are not an isolated real-time deadline benchmark. Positive clearance below
the node margin is recorded without reclassifying it as collision.

\section{Reproduction and artifacts}
Run \texttt{runtests('tests')} from the repository root. The paired study is
run using the following entry point:
\begin{quote}\small
\path{scripts/verifyExperimentalAssumptionPolicy.m}
\end{quote}
Supply an external output directory and a 1200-hold cap. Compact final test counts, scenario
results and source provenance are in
\path{report/EXPERIMENTAL_ASSUMPTION_POLICY_20261003/}.
Full traces, continuation files, console output and this report's derived PDF
are retained as identified external archive exports. Operational source and
tests are committed normally. Unrelated working-tree changes and external
solver dependencies are excluded.
\end{document}
